% from ega2-1-fr.tex
\oldpage[II]{197}
  On dit que $h(x,0)$ est déduit de $x$ par \emph{homogénéisation} à l'aide
  de la «~variable homogénéisante~» $z$. Mais comme $h(x,0)$ ne dépend pas
  additivement de $x$ (ni \emph{a fortiori} $S$-linéairement), \emph{on se
  gardera de croire} (même lorsque $M=S$) que les $h(x,0)$ parcourent un
  \emph{système de générateurs} du $\widehat{S}$-module gradué
  $\overline{N}$ lorsqu'on fait parcourir à $x$ un \emph{système de
  générateurs} du $S$-module $N$. Il en est cependant bien ainsi dans le cas
  (seul considéré dans la géométrie algébrique élémentaire) où $N$ est un
  $S$-module \emph{libre monogène}, car si $t$ est une base de $N$, $h(t,0)$
  engendre le $\widehat{S}$-module $\overline{N}$.
\end{enumerate}
\end{remark}

\subsection{Compléments sur les faisceaux associés aux
$\mathcal{S}$-Modules gradués}
\label{subsection:II.8.14-fr}

\begin{env}[8.14.1]
\label{II.8.14.1-fr}
Soient $Y$ un préschéma, $\mathcal{S}$ une $\mathcal{O}_Y$-Algèbre graduée
quasi-cohérente à \emph{degrés positifs},
$X=\operatorname{Proj}(\mathcal{S})$, $q:X\to Y$ le morphisme structural
(séparé en vertu de (\hyperref[II.3.1.3-fr]{3.1.3})). Conservons les
notations de (\hyperref[II.8.12.1-fr]{8.12.1}), de sorte que nous avons
défini un foncteur
$\mathcal{M}_X=\operatorname{Proj}(\mathcal{M})$ en $\mathcal{M}$, de la
catégorie des $\mathcal{S}$-Modules gradués quasi-cohérents dans celle des
$\mathcal{S}_X$-Modules gradués quasi-cohérents~; il est clair en outre
(\hyperref[II.3.2.4-fr]{3.2.4}) que c'est un foncteur \emph{additif exact},
commutant aux limites inductives.

Notons en outre qu'il résulte aussitôt de la définition
(\hyperref[II.8.12.1.1-fr]{8.12.1.1}) que l'on a
\begin{equation}
\label{II.8.14.1.1-fr}
  \operatorname{Proj}(\mathcal{M}(n))
  =(\operatorname{Proj}(\mathcal{M}))(n)
  \qquad\text{pour tout }n\in\mathbf{Z}.
\tag{8.14.1.1}
\end{equation}
\end{env}

\begin{env}[8.14.2]
\label{II.8.14.2-fr}
Nous allons d'abord étendre aux $\mathcal{S}_X$-Modules de la forme
$\operatorname{Proj}(\mathcal{M})$ les homomorphismes canoniques $\lambda$
et $\mu$ définis dans (\hyperref[II.3.2.6-fr]{3.2.6}). Pour cela, notons
que, quels que soient $m\in\mathbf{Z}$ et $n\in\mathbf{Z}$, on a, en vertu
de (\hyperref[II.2.1.2.1-fr]{2.1.2.1}), un homomorphisme canonique de
$\mathcal{O}_X$-Modules
\begin{equation}
\label{II.8.14.2.1-fr}
  \lambda_{mn}:\operatorname{Proj}_0(\mathcal{M}(m))
  \otimes_{\mathcal{O}_X}\operatorname{Proj}_0(\mathcal{N}(n))
  \to\operatorname{Proj}_0
  ((\mathcal{M}\otimes_{\mathcal{S}}\mathcal{N})(m+n))
\tag{8.14.2.1}
\end{equation}
et de même, en vertu de (\hyperref[II.2.1.2.2-fr]{2.1.2.2}), un
homomorphisme canonique de $\mathcal{O}_X$-Modules
\begin{equation}
\label{II.8.14.2.2-fr}
  \mu_{mn}:\operatorname{Proj}_0
  ((\mathcal{H}om_{\mathcal{S}}(\mathcal{M},\mathcal{N}))(n-m))
  \to\mathcal{H}om_{\mathcal{O}_X}
  (\operatorname{Proj}_0(\mathcal{M}(m)),
   \operatorname{Proj}_0(\mathcal{N}(n)))
\tag{8.14.2.2}
\end{equation}
quels que soient les $\mathcal{S}$-Modules gradués quasi-cohérents
$\mathcal{M}$, $\mathcal{N}$. On en déduit un homomorphisme
\[
  \mu_k:\operatorname{Proj}_0
  ((\mathcal{H}om_{\mathcal{S}}(\mathcal{M},\mathcal{N}))(k))
  \to
  (\mathcal{H}om_{\mathcal{S}_X}
  (\operatorname{Proj}(\mathcal{M}),\operatorname{Proj}(\mathcal{N})))_k
\]
obtenu en faisant correspondre à tout
$u\in\Gamma(U,\operatorname{Proj}_0
((\mathcal{H}om_{\mathcal{S}}(\mathcal{M},\mathcal{N}))(k)))$
l'homomorphisme $\mu_k(u)$, de degré $k$, de $\mathbf{Z}$-modules gradués
$\Gamma(U,\operatorname{Proj}(\mathcal{M}))\to
\Gamma(U,\operatorname{Proj}(\mathcal{N}))$ ($U$ ouvert de $X$) qui dans
chaque $\Gamma(U,\operatorname{Proj}_0(\mathcal{M}(m)))$ coïncide avec
$\mu_{m,m+k}(u)$~; en outre, en revenant à la définition des $\mu_{mn}$
(\hyperref[II.2.5.12.1-fr]{2.5.12.1}), on voit aussitôt que $\mu_k(u)$ est
en fait un homomorphisme de degré $k$ de
$\Gamma(U,\mathcal{S}_X)$-modules gradués, et en outre que les $\mu_k$
définissent un homomorphisme de $\mathcal{S}_X$-Modules gradués
\begin{equation}
\label{II.8.14.2.3-fr}
  \operatorname{Proj}
  (\mathcal{H}om_{\mathcal{S}}(\mathcal{M},\mathcal{N}))
  \to\mathcal{H}om_{\mathcal{S}_X}
  (\operatorname{Proj}(\mathcal{M}),\operatorname{Proj}(\mathcal{N})).
\tag{8.14.2.3}
\end{equation}

De même, compte tenu du diagramme d'associativité
(\hyperref[II.2.5.11.4-fr]{2.5.11.4}), les homomorphismes
(\hyperref[II.8.14.2.1-fr]{8.14.2.1}) donnent un homomorphisme de
$\mathcal{S}_X$-Modules gradués
\begin{equation}
\label{II.8.14.2.4-fr}
  \lambda:\operatorname{Proj}(\mathcal{M})
  \otimes_{\mathcal{S}_X}\operatorname{Proj}(\mathcal{N})
  \to\operatorname{Proj}(\mathcal{M}\otimes_{\mathcal{S}}\mathcal{N}).
\tag{8.14.2.4}
\end{equation}
\end{env}

